\subsection{Fundamental Set Calculations}

Set calculations combine or compare membership domains. These operations do not care about positions, because a set has no positional structure. They care about membership: which elements are present in the left set, which elements are present in the right set, and which elements should appear in the result.

The four fundamental calculations are:

\begin{itemize}
    \item union: everything from both sets;
    \item intersection: only what appears in both sets;
    \item difference: what appears in the first set but not the second;
    \item symmetric difference: what appears in exactly one of the two sets.
\end{itemize}

Each operation has an operator form and a method form. The operator form is compact. The method form is more explicit and can accept general iterables in more cases.

\subsubsection{The Union Operator \texttt{|} and Method Equivalent \texttt{.union()}}

Union builds a new set containing every distinct element from both input sets.

\begin{verbatim}
python_lines = {"spam", "eggs", "larch"}
tv_lines = {"D'oh!", "spam", "Serenity now!"}

result1 = python_lines | tv_lines
result2 = python_lines.union(tv_lines)

print(f"python_lines: {python_lines!r}")
print(f"tv_lines:     {tv_lines!r}")
print()
print(f"result1: {result1!r}")
print(f"result2: {result2!r}")
print(f"same result: {result1 == result2}")
print(f"type(result1): {type(result1)}")
\end{verbatim}

Possible output:

\begin{verbatim}
python_lines: {'eggs', 'larch', 'spam'}
tv_lines:     {"D'oh!", 'Serenity now!', 'spam'}

result1: {'Serenity now!', 'eggs', "D'oh!", 'larch', 'spam'}
result2: {'Serenity now!', 'eggs', "D'oh!", 'larch', 'spam'}
same result: True
type(result1): <class 'set'>
\end{verbatim}

The value \texttt{"spam"} appears in both input sets, but it appears only once in the result. Union preserves the set uniqueness rule.

The original sets are not changed.

\begin{verbatim}
left = {"Arthur", "Brian"}
right = {"Brian", "Lancelot"}

result = left | right

print(f"left:   {left!r}")
print(f"right:  {right!r}")
print(f"result: {result!r}")
\end{verbatim}

Possible output:

\begin{verbatim}
left:   {'Arthur', 'Brian'}
right:  {'Lancelot', 'Brian'}
result: {'Arthur', 'Lancelot', 'Brian'}
\end{verbatim}

The operator \texttt{|} expects set-like operands. The method \texttt{.union()} can accept other iterable objects.

\begin{verbatim}
base = {"Homer", "Marge"}
more_names = ["Lisa", "Bart", "Homer"]

result = base.union(more_names)

print(f"base: {base!r}")
print(f"more_names: {more_names!r}")
print(f"result: {result!r}")
\end{verbatim}

Possible output:

\begin{verbatim}
base: {'Homer', 'Marge'}
more_names: ['Lisa', 'Bart', 'Homer']
result: {'Lisa', 'Homer', 'Bart', 'Marge'}
\end{verbatim}

Here \texttt{more\_names} is a list, not a set. The method reads its elements and inserts them into the resulting set.

\subsubsection{The Intersection Operator \texttt{\&} and Method Equivalent \texttt{.intersection()}}

Intersection builds a new set containing only elements that appear in both input sets.

\begin{verbatim}
python_lines = {"spam", "eggs", "larch", "ni"}
famous_words = {"spam", "D'oh!", "ni", "yada yada"}

result1 = python_lines & famous_words
result2 = python_lines.intersection(famous_words)

print(f"python_lines: {python_lines!r}")
print(f"famous_words: {famous_words!r}")
print()
print(f"result1: {result1!r}")
print(f"result2: {result2!r}")
print(f"same result: {result1 == result2}")
\end{verbatim}

Possible output:

\begin{verbatim}
python_lines: {'eggs', 'larch', 'spam', 'ni'}
famous_words: {'yada yada', "D'oh!", 'spam', 'ni'}

result1: {'spam', 'ni'}
result2: {'spam', 'ni'}
same result: True
\end{verbatim}

The result contains \texttt{"spam"} and \texttt{"ni"} because those values appear in both sets.

Intersection is useful when two independent domains must both approve the same value.

\begin{verbatim}
allowed_users = {"Jerry", "Elaine", "George", "Kramer"}
present_users = {"Elaine", "Newman", "Kramer"}

active_allowed = allowed_users & present_users

print(f"allowed_users:  {allowed_users!r}")
print(f"present_users:  {present_users!r}")
print(f"active_allowed: {active_allowed!r}")
\end{verbatim}

Possible output:

\begin{verbatim}
allowed_users:  {'George', 'Elaine', 'Jerry', 'Kramer'}
present_users:  {'Elaine', 'Newman', 'Kramer'}
active_allowed: {'Elaine', 'Kramer'}
\end{verbatim}

The method form can also accept a general iterable.

\begin{verbatim}
known = {"Homer", "Marge", "Lisa", "Bart"}
incoming = ["Lisa", "Maggie", "Bart", "Burns"]

result = known.intersection(incoming)

print(f"known: {known!r}")
print(f"incoming: {incoming!r}")
print(f"result: {result!r}")
\end{verbatim}

Possible output:

\begin{verbatim}
known: {'Lisa', 'Homer', 'Bart', 'Marge'}
incoming: ['Lisa', 'Maggie', 'Bart', 'Burns']
result: {'Lisa', 'Bart'}
\end{verbatim}

Again, the result is a set. The list is only used as a source of candidate values.

\subsubsection{The Difference Operator \texttt{-} and Method Equivalent \texttt{.difference()}}

Difference builds a new set containing elements from the left set that do not appear in the right set.

\begin{verbatim}
all_items = {"spam", "eggs", "larch", "beans"}
blocked_items = {"eggs", "beans"}

result1 = all_items - blocked_items
result2 = all_items.difference(blocked_items)

print(f"all_items:     {all_items!r}")
print(f"blocked_items: {blocked_items!r}")
print()
print(f"result1: {result1!r}")
print(f"result2: {result2!r}")
print(f"same result: {result1 == result2}")
\end{verbatim}

Possible output:

\begin{verbatim}
all_items:     {'eggs', 'beans', 'larch', 'spam'}
blocked_items: {'eggs', 'beans'}

result1: {'larch', 'spam'}
result2: {'larch', 'spam'}
same result: True
\end{verbatim}

The result keeps only the elements that belong to \texttt{all\_items} and do not belong to \texttt{blocked\_items}.

Difference is directional. Reversing the operands changes the question.

\begin{verbatim}
a = {"Arthur", "Brian", "Lancelot"}
b = {"Brian", "Galahad"}

print(f"a - b: {a - b}")
print(f"b - a: {b - a}")
\end{verbatim}

Possible output:

\begin{verbatim}
a - b: {'Arthur', 'Lancelot'}
b - a: {'Galahad'}
\end{verbatim}

The first expression asks: what is in \texttt{a} but not in \texttt{b}? The second expression asks: what is in \texttt{b} but not in \texttt{a}? These are not the same question.

The method form can accept an iterable.

\begin{verbatim}
seen = {"Homer", "Marge", "Lisa", "Bart"}
already_processed = ["Homer", "Lisa"]

remaining = seen.difference(already_processed)

print(f"seen: {seen!r}")
print(f"already_processed: {already_processed!r}")
print(f"remaining: {remaining!r}")
\end{verbatim}

Possible output:

\begin{verbatim}
seen: {'Lisa', 'Homer', 'Bart', 'Marge'}
already_processed: ['Homer', 'Lisa']
remaining: {'Bart', 'Marge'}
\end{verbatim}

The original set is not changed. A new set is produced.

\subsubsection{The Symmetric Difference Operator \texttt{\textasciicircum{}} and Method Equivalent \texttt{.symmetric\_difference()}}

Symmetric difference builds a new set containing elements that appear in exactly one of the two sets. It excludes elements that appear in both.

\begin{verbatim}
left = {"spam", "eggs", "larch"}
right = {"eggs", "beans", "larch", 42}

result1 = left ^ right
result2 = left.symmetric_difference(right)

print(f"left:  {left!r}")
print(f"right: {right!r}")
print()
print(f"result1: {result1!r}")
print(f"result2: {result2!r}")
print(f"same result: {result1 == result2}")
\end{verbatim}

Possible output:

\begin{verbatim}
left:  {'eggs', 'larch', 'spam'}
right: {42, 'eggs', 'beans', 'larch'}

result1: {42, 'beans', 'spam'}
result2: {42, 'beans', 'spam'}
same result: True
\end{verbatim}

The values \texttt{"eggs"} and \texttt{"larch"} are excluded because they appear in both sets. The values \texttt{"spam"}, \texttt{"beans"}, and \texttt{42} remain because each appears on only one side.

Symmetric difference is useful when the program needs disagreement between two domains.

\begin{verbatim}
expected = {"Jerry", "Elaine", "George"}
actual = {"Jerry", "Elaine", "Kramer"}

mismatch = expected ^ actual

print(f"expected: {expected!r}")
print(f"actual:   {actual!r}")
print(f"mismatch: {mismatch!r}")
\end{verbatim}

Possible output:

\begin{verbatim}
expected: {'Elaine', 'George', 'Jerry'}
actual:   {'Elaine', 'Jerry', 'Kramer'}
mismatch: {'George', 'Kramer'}
\end{verbatim}

The mismatch contains \texttt{"George"} because he was expected but not actual, and \texttt{"Kramer"} because he was actual but not expected.

The method form can accept an iterable.

\begin{verbatim}
base = {"Arthur", "Brian", "Lancelot"}
other = ["Brian", "Galahad", "Arthur"]

result = base.symmetric_difference(other)

print(f"base: {base!r}")
print(f"other: {other!r}")
print(f"result: {result!r}")
\end{verbatim}

Possible output:

\begin{verbatim}
base: {'Arthur', 'Brian', 'Lancelot'}
other: ['Brian', 'Galahad', 'Arthur']
result: {'Galahad', 'Lancelot'}
\end{verbatim}

The values \texttt{"Arthur"} and \texttt{"Brian"} are removed from the result because they appear in both sources. The values \texttt{"Lancelot"} and \texttt{"Galahad"} remain because each appears only on one side.

The four calculations can be summarized compactly:

\begin{center}
\begin{tabular}{llll}
\textbf{Calculation} & \textbf{Operator} & \textbf{Method} & \textbf{Meaning} \\
\hline
Union & \texttt{|} & \texttt{.union()} & In either set \\
Intersection & \texttt{\&} & \texttt{.intersection()} & In both sets \\
Difference & \texttt{-} & \texttt{.difference()} & In left, not in right \\
Symmetric difference & \texttt{\textasciicircum{}} & \texttt{.symmetric\_difference()} & In exactly one set \\
\end{tabular}
\end{center}

All four operations produce new sets. They do not mutate the left operand.

\begin{verbatim}
a = {"spam", "eggs"}
b = {"eggs", "larch"}

union_result = a | b
intersection_result = a & b
difference_result = a - b
symmetric_result = a ^ b

print(f"a after calculations: {a!r}")
print(f"b after calculations: {b!r}")
print()
print(f"union_result:        {union_result!r}")
print(f"intersection_result: {intersection_result!r}")
print(f"difference_result:   {difference_result!r}")
print(f"symmetric_result:    {symmetric_result!r}")
\end{verbatim}

Possible output:

\begin{verbatim}
a after calculations: {'eggs', 'spam'}
b after calculations: {'eggs', 'larch'}

union_result:        {'eggs', 'larch', 'spam'}
intersection_result: {'eggs'}
difference_result:   {'spam'}
symmetric_result:    {'larch', 'spam'}
\end{verbatim}